\chapter{\MakeUppercase{Technology / Paper Overview}}
\setcounter{equation}{0}

\section{Selected Paper}
J. Sun et al., "Multimodal Behavioral Data Predict Stroke Patient’s Response to BCI Treatment Through Explainable AI," in IEEE Transactions on Neural Systems and Rehabilitation Engineering, vol. 34, pp. 1-9, 2026, doi: 10.1109/TNSRE.2025.3625222.\cite{bib1}

\section{Objectives}
\begin{itemize}
    \item Investigate the relationship between baseline characteristics and post-treatment motor function.
    \item Determine whether combining demographic, clinical and biomechanical features improves prediction.
    \item Classify patients as responders or non-responders to BCI rehabilitation.
    \item Compare Elastic Net and ANN for treatment-response classification.
    \item Identify important predictors using stepwise regression and SHAP.
    \item Provide interpretable evidence that may support personalized rehabilitation planning.
\end{itemize}

\section{Architecture}
The proposed framework uses multimodal behavioral data collected before BCI rehabilitation to predict patients’ motor recovery outcomes after two weeks of treatment. The input data include demographic information, clinical assessment scores, and biomechanical measurements. These features are processed using regression and classification models, including Elastic Net and Artificial Neural Networks (ANN). Explainable AI techniques, particularly SHAP, are then used to identify the features that contribute to the predictions. The overall workflow and detailed methodology are presented in the next chapter. [1]
\section{Key Contributions of the Authors}

\begin{itemize}
    \item Combines demographic, clinical, and biomechanical data to predict rehabilitation outcomes.
    \item Evaluates prediction models for estimating post-treatment lower-limb motor function.
    \item Classifies patients according to their response to BCI rehabilitation.
    \item Compares the predictive performance of Elastic Net and ANN.
    \item Uses explainable AI techniques to identify important recovery predictors.
    \item Highlights the importance of baseline motor function, muscle spasticity, and balance ability in predicting rehabilitation outcomes \cite{bib1}.
\end{itemize}


